\section{Discussion}
\label{sec:discussion}

The central result is a decision consequence: in this benchmark, 269 of 277 cells that pass the conventional error-based rule do not pass the added fidelity requirements. The audit therefore exposes a difference between being better than persistence under RMSE and retaining the observed variance of the forecast trajectory.

The discordant cells clarify the practical meaning of that difference. A model can pass the error and P75-recall screens while producing a forecast with $\alpha<0.50$. Conversely, seasonal naive retains variance but is rarely error-eligible, while STL+Ridge avoids collapse at the cost of poor error skill. These profiles show why error and fidelity should be reported together when model selection is intended to support dynamic episode decisions.

The Rule-B thresholds are deliberately operational. Their purpose here is sensitivity of the eligibility decision, not prescription of a universal standard. The appropriate fidelity threshold depends on the decision, pollutant, lead time, and tolerance for missed episodes or false alarms. The result supports an auditable multi-dimensional reporting practice, not a new metric or a replacement for RMSE.

\subsection{Limitations}
The evidence is bounded to daily PM$_{10}$ from 17 Spanish monitoring stations, five model families, seven horizons, and a rolling-origin benchmark. The 595 cells share stations and underlying series and are not independent replicates. Rule-B thresholds are study-specific. The aggregate evidence is recovered and frozen; the row-level manifest notes that SARIMA skill is not independently reconstructable cell-by-cell from the distributed combined row-level file alone. The analysis uses deterministic point forecasts, does not establish causal explanations for variance collapse, and does not establish generalisation to other pollutants, countries, temporal resolutions, horizons, or model families.
